\documentclass[a4paper,12pt]{article}
\usepackage{color}
\usepackage{graphicx, gensymb}
\usepackage{amsfonts, cancel, amssymb}
\usepackage{amsmath, amsthm, array, esvect, esint}
\usepackage{typearea, tikz, pgffor, verbatim}
\usetikzlibrary{calc, quotes}
\usepackage[a4paper,left=2cm,right=2cm,top=2cm,bottom=3cm,bindingoffset=5mm]{geometry}
\newcommand\numberthis{\addtocounter{equation}{1}\tag{\theequation}}
\newcommand{\bra}{}
\newcommand{\kom}{}
\newcommand{\brak}{}
\newcommand{\braket}{}
\newcommand{\intx}{}
\newcommand{\intr}{}
\newcommand{\kla}{}
\newcommand{\klam}{}
\newcommand{\klammer}{}
\newcommand{\oper}{}
\newcommand{\tecxt}{}
\newcommand{\Bra}{}
\newcommand{\Ket}{}

\begin{document}

\title{02 Informationstheorie}
\author{Joachim Schwardt}
\date{\today}
\maketitle

\section*{Aufgabe 2.1: Entropie}
Die Entropie sei definiert als $S=-\sum_{i\in\Omega}p(i)\log_2[p(i)]$. Setze $p(0)=p_0$ und $p(1)=p_1$, wobei natürlich $p_0+p_1=1$ gilt. 
\subsection*{a)}Hier ist $\Omega=\{0,1\}$ mit $|\Omega|=2$:
\begin{align*}
S_a&=-p_0\log_2p_0-p_1\log_2p_1\kom\overset{\substack{{p_0=1/2} \\ |}}{=}1\numberthis\label{S aus a}
\end{align*}
\subsection*{b)}Hier ist $\Omega=\{\{0,0\},\{0,1\},\{1,0\},\{1,1\}\}$ mit $|\Omega|=4$:
\begin{align*}
S_b&=-p_0^2\log_2p_0^2-2p_0p_1\log_2p_0p_1-p_1^2\log_2p_1^2=\\
&=-(p_0+p_1)2p_0\log_2p_0-(p_0+p_1)2p_1\log_2p_1=2S_a\kom\overset{\substack{{p_0=1/2} \\ |}}{=}2\numberthis\label{S aus b}
\end{align*}
\subsection*{c)}Hier ist $\Omega=\{\{0,0\},\{0,1\},\{1,1\}\}$ mit $|\Omega|=3$:
\begin{align*}
S_c&=-p_0^2\log_2p_0^2-2p_0p_1\log_22p_0p_1-p_1^2\log_2p_1^2=\\
&=S_b-2p_0p_1\log_22=S_b-2p_0p_1\kom\overset{\substack{{p_0=1/2} \\ |}}{=}\frac{3}{2}\numberthis\label{S aus c}
\end{align*}
\subsection*{d)}Hier ist $\Omega=\{\{0,\dots,0,0\},\{0,\dots,0,1\},\{0,\dots,1,0\},\dots,\{1,\dots,1,1\}\}^n$ mit $|\Omega|=2^n$:
\begin{align*}
S_d&=-\sum_{k=0}^n\binom{n}{k}p_0^kp_1^{n-k}\log_2p_0^kp_1^{n-k}=\\
&=-\log_2p_0\sum_{k=0}^n\binom{n}{k}kp_0^kp_1^{n-k}-\log_2p_1\sum_{k=0}^n\binom{n}{k}(n-k)p_0^kp_1^{n-k}=\\
&=-np_0\log_2p_0-np_1\log_2p_1=nS_a\kom\overset{\substack{{p_0=1/2} \\ |}}{=}n\numberthis\label{S aus d}
\end{align*}
\subsection*{e)}Hier ist $\Omega=\{\{0,\dots,0,0\},\{0,\dots,0,1\},\{0,\dots,1,1\},\dots,\{1,\dots,1,1\}\}$ mit $|\Omega|=n+1$:
\begin{align*}
S_e&=-\sum_{k=0}^n\binom{n}{k}p_0^kp_1^{n-k}\log_2\binom{n}{k}p_0^kp_1^{n-k}=S_d-\sum_{k=0}^n\binom{n}{k}p_0^kp_1^{n-k}\log_2\binom{n}{k}\kom\overset{\substack{{p_0=1/2} \\ |}}{=}\\
&=n-\frac{1}{2^n}\sum_{k=0}^n\binom{n}{k}\log_2\binom{n}{k}\numberthis\label{S aus e}
\end{align*}
\subsection*{f)}
Sei $p(X_{1,2})=p_{1,2}$ und $X_{1,2}$ unabhängige Ereignisse. Dann gilt $p(X_1,X_2)=p_1(X_1)p_2(X_2)$, woraus für die Entropie folgt:
\begin{align*}
S_{1,2}&=-\sum_{ij}p_1(i)p_2(j)\log_2p_1(i)p_2(j)=\\
&=-\sum_ip_1(i)\log_2p_1(i)\sum_jp_2(j)-\sum_jp_2(j)\log_2p_2(j)\sum_ip_1(i)=S_1+S_2
\end{align*}
\section*{Aufgabe 2.2: Entropie einer unfairen Münze}
\subsection*{a)}Sei $p(0)=p$. Dann erwarten wir folgende Anzahl an Bits:
\begin{align*}
N&=p^2\cdot 1+p(1-p)\cdot 2+(1-p)p\cdot 3+(1-p)^2\cdot 3=3-p-p^2
\end{align*}
Um im Mittel weniger als 2 Bits zu benötigen, fordern wir also $N<2$. Das führt auf die Bedingung $p>\frac{\sqrt{5}-1}{2}=\frac{1}{\phi}\approx0.618$.
\subsection*{b)}
Bei $n$ Münzwürfen ergibt sich die mittlere Anzahl $n_0$ an Nullen gemäß
\begin{align*}
n_0=\sum_{k=0}^nkp^k(1-p)^{n-k}=np
\end{align*} 
Die Standardabweichung erhält man aus $\sigma=\sqrt{\tecxt\text{Var}}$
\begin{align*}
\tecxt\text{Var}&=\sum_{k=0}^n(k-n_0)^2k^p(1-p)^{n-k}=n_0p(1-p)
\end{align*}
Es gibt $\binom{n}{n_0}$ verschiedene Konfigurationen mit genau $n_0$ Nullen. Für große $n$ steigt der Rechenaufwand allerdings enorm, weshalb die Stirling-Näherung verwendet werden soll:
\begin{align*}
N(p)&=\frac{n!}{np!(n-np)!}\approx\frac{(\frac{n}{e})^n}{\sqrt{2\pi p(n-np)}(\frac{np}{e})^{np}(\frac{n-np}{e})^{n-np}}=\frac{p^{-np}(1-p)^{-n(1-p)}}{\sqrt{2\pi np(1-p)}}\numberthis\label{equation:N von p}
\end{align*}
Damit folgt für das Produkt $\sigma N(p)$
\begin{align*}
\sigma N&\approx \frac{1}{\sqrt{2\pi}}p^{-np}(1-p)^{-n(1-p)}\numberthis\label{equation: sigma N von p}
\end{align*}
Zuletzt betrachten wird das Verhältnis $\alpha:=\frac{\sigma N}{2^{nS(p)}}$:
\begin{align*}
\alpha &=\sigma N \klam\left[ {2^{p\log_2p+(1-p)\log_2(1-p)}} \right]^n=\sigma Np^{np}(1-p)^{n(1-p)}=\frac{1}{\sqrt{2\pi}}
\end{align*}
\section*{Aufgabe 2.3: Kombinatorik, Entropie und Sudoku}
\subsection*{a)}Jede Konfiguration hat die gleiche Wahrscheinlichkeit $p(i)=\frac{1}{N_0}$. Damit folgt für die Entropie (zur Basis 9):
\begin{align*}
S_a&=-\sum_ip(i)\log_9p(i)=\log_9N_0\approx 22,87
\end{align*}
\subsection*{b)}Ich rate $S=1$, da ja genau ein Bit (also ein Bit mit neun möglichen Konfigurationen; Logarithmus zur Basis 9) benötigt wird, um die Information zu speichern.
\subsection*{c)}
Durch die Regeln von Sudoku ist es egal, ob in einer Reihe 8 oder 9 Zahlen vorgegeben werden, da sich die "fehlende" einfach bestimmen lässt.
\subsection*{d)}
Die Entropie ist in beiden Fällen einfach 0, da die Lösung hier eindeutig ist. Folglich gibt es überhaupt keine fehlende Information mehr, womit es auch keine Entropie gibt. 
\section*{Aufgabe 2.4: Ideales Gas, thermisches Gleichgewicht}
Die Ideale Gasgleichung ist $pV=k_BNT$. Wir betrachten zwei Systeme mit gegebenen $p_i,T_i,V_i$, die allerdings das gleiche ideale Gas enthalten.
\subsection*{a)}Wärmeaustausch mit einem Wärmebad $T_0$, also gilt $T_i'=T_0$. Die Systeme können nach Voraussetzung allerdings nur Wärme tauschen, weshalb hier nur noch die Interaktion mit dem Wärmebad relevant ist. Es genügt also, nur ein System zu betrachten.
\begin{align*}
\frac{p'V'}{pV}&=\frac{T_0}{T}
\end{align*}
Ohne einen weiteren Zusammenhang zwischen $p$ und $V$ zu kennen, kann aber nichts weiter bestimmt werden. Es gilt für die Änderung der Temperatur einfach $\Delta U_i=k_B(T_0-T_i)$.
\subsection*{b)}
Hier haben beide Systeme bereits die gleiche Temperatur $T_1=T_2$, sie können aber Volumen austauschen. Damit erhält das System einen weiteren Freiheitsgrad, denn die einzige Gleichung enthält nun drei Unbekannt:
\begin{align*}
\frac{p'V'}{pV}&=\frac{N'}{N}
\end{align*}
\subsection*{c)}

\end{document}